\chapter{Thinking in Expected Value}\label{ch:m2:thinking-in-expected-value}

In 2016 two researchers gave 61 young finance professionals and students 25 dollars
and 30 minutes to bet on a coin that, they were told, landed heads 60\% of the time. They
could bet any amount on either side, as often as they liked, up to an undisclosed
maximum payout of 250 dollars. A fixed bet of 10 to 20\% of the bankroll on heads would
have reached the maximum 95\% of the time. Only 21\% of the players reached it; a third
ended with less than they started with, and 28\% went bust, losing everything on a coin
tilted in their favour. They bet too much, bet erratically, and some bet on tails. This
chapter is about the arithmetic they skipped: edge and \omterm{def:m2:thinking-in-expected-value:edge}{odds}, the fraction of wealth that
maximises long-run growth, why traders stake less than that, and how much less when the
edge itself is only an estimate.

\section{Edge and odds}

\begin{definition}[Edge, odds]\label{def:m2:thinking-in-expected-value:edge}
The \emph{odds}\index{odds} $b$ of a bet are what it pays per unit staked if it wins,
the stake being lost otherwise. Its \emph{edge}\index{edge} is its expected profit per
unit staked, $pb - q$, where $p$ is the probability of winning and $q = 1 - p$.
\end{definition}

The coin pays even money, $b = 1$, and wins with $p = 0.6$: its edge is $0.6 - 0.4 = 0.2$,
20 cents per dollar staked. A positive edge says the bet is worth taking; it says nothing
about how much to stake. Stake everything each time and a single tails ends the game:
the expected value of each bet is positive, but the probability of surviving 300 flips
is $0.6^{300}$. The question is not the edge of one bet but what staking a given fraction
of wealth, bet after bet, does to wealth.

\begin{definition}[Growth rate]\label{def:m2:thinking-in-expected-value:growth}
The \emph{growth rate}\index{growth rate} of a staking rule is the expected logarithm of
the ratio of wealth after a bet to wealth before it. Staking a fraction $f$ of wealth on
a bet with \omterm{def:m2:thinking-in-expected-value:edge}{odds} $b$ and win probability $p$, it is
\[
g(f) = p\ln(1 + fb) + q\ln(1 - f).
\]
\end{definition}

After $n$ bets, the logarithm of wealth is a sum of $n$ independent terms, so by the law
of large numbers wealth grows almost surely like $e^{n g(f)}$: the \omterm{def:m2:thinking-in-expected-value:growth}{growth rate}, not the
edge, decides where wealth ends up. The expected wealth after $n$ bets is $(1 + f(pb -
q))^n$, which grows with $f$ without limit; the median, $e^{n g(f)}$ at most, does not. At
$f = 0.2$ on the coin, the expected gain is 4\% a flip, and the expected value of 300 flips
from 25 dollars is USD~3.2 million, the figure Haghani and Dewey quote; the median is
USD~10\,504.

\section{The Kelly criterion}

\begin{definition}[Kelly criterion]\label{def:m2:thinking-in-expected-value:kelly}
The \emph{Kelly criterion}\index{Kelly criterion} stakes, on each bet, the fraction of
wealth that maximises the \omterm{def:m2:thinking-in-expected-value:growth}{growth rate}.
\end{definition}

\begin{proposition}[The Kelly fraction]\label{prop:m2:thinking-in-expected-value:kelly}
For a bet with \omterm{def:m2:thinking-in-expected-value:edge}{odds} $b$ and win probability $p$, the \omterm{def:m2:thinking-in-expected-value:growth}{growth rate} is maximised at
\[
f^* = p - \frac{q}{b} = \frac{\text{edge}}{\text{odds}},
\]
if the edge is positive, and at $f^* = 0$ otherwise. For small bets with mean $\mu$ and
variance $s^2$ per unit staked, $g(f) \approx f\mu - \tfrac12 f^2 s^2$, maximised at $f^* =
\mu / s^2$, with $g(f^*) = \mu^2 / (2s^2)$.
\end{proposition}

\begin{proof}
$g'(f) = pb/(1 + fb) - q/(1 - f)$ vanishes at $f = (pb - q)/b$, and $g$ is concave. For the
second form, expand the logarithm to second order: $\ln(1 + fX) \approx fX - \tfrac12 f^2X^2$
and take expectations.
\end{proof}

The coin's Kelly fraction is 0.2: bet 20\% of the bankroll on heads each time, for a \omterm{def:m2:thinking-in-expected-value:growth}{growth
rate} of 2.01\% a flip. The quadratic form shows the shape of the trade-off
(\cref{fig:m2:thinking-in-expected-value:growth}): growth rises, peaks at $f^*$, and falls
back to zero at about twice $f^*$; beyond that, the more one bets, the faster one is ruined,
however large the edge. On the coin the \omterm{def:m2:thinking-in-expected-value:growth}{growth rate} is zero at a stake of 38.9\%.

\begin{omfigure}
\begin{tikzpicture}
\begin{axis}[width=11cm,height=4.8cm,
  xlabel={fraction of wealth staked each flip, $f$},ylabel={growth rate (\% per flip)},
  tick label style={font=\small},label style={font=\small},
  xmin=0,xmax=0.46,ymin=-1,ymax=2.8,grid=major,grid style={gray!25},
  xticklabel style={/pgf/number format/.cd,fixed,precision=2},
  table/col sep=comma]
\addplot[omDef,very thick] table[x=f,y=g]{figdata/markets-2/29-thinking-in-expected-value/growth.csv};
\addplot[gray,thin] coordinates {(0,0) (0.46,0)};
\addplot[only marks,mark=*,mark size=1.8pt,black] coordinates {(0.2,2.0136) (0.3894,0)};
\node[font=\footnotesize,anchor=south] at (axis cs:0.2,2.05) {Kelly: 0.2};
\node[font=\footnotesize,anchor=east] at (axis cs:0.375,-0.45) {zero growth: 0.389};
\end{axis}
\end{tikzpicture}

\omcaption{\omterm{def:m2:thinking-in-expected-value:growth}{Growth rate} of wealth per flip of the 60\% coin at even money, against the
fraction of wealth staked. It peaks at the Kelly fraction, 0.2, with 2.01\% a flip, and
falls to zero at 0.389, about twice Kelly. Data: the chapter's tutorial.}\label{fig:m2:thinking-in-expected-value:growth}
\end{omfigure}

\section{Fractional Kelly and risk of ruin}

\begin{definition}[Fractional Kelly, drawdown, risk of ruin]\label{def:m2:thinking-in-expected-value:frac}
\emph{Fractional Kelly}\index{fractional Kelly} stakes a multiple $c < 1$ of the Kelly
fraction. A \emph{drawdown}\index{drawdown} is a fall of wealth from its previous peak,
measured as a fraction of that peak. The \emph{risk of ruin}\index{risk of ruin} of a
staking rule is the probability that wealth ever falls to a given fraction of its
starting value.
\end{definition}

Kelly staking maximises growth, but it is a violent ride: the distribution of wealth
after 300 flips is enormously spread (\cref{fig:m2:thinking-in-expected-value:coin}).
Half Kelly gives up a quarter of the \omterm{def:m2:thinking-in-expected-value:growth}{growth rate} in the quadratic approximation, $g(cf^*)
= (2c - c^2)g(f^*)$, and much of the spread: on the coin, the median after 300 flips is
USD~2\,279 against 10\,504, but the worst tenth of outcomes is better, 251 against 121.
Twice Kelly has, to second order, zero growth: half its players end below their stake.

\begin{omfigure}
\begin{tikzpicture}
\begin{axis}[width=11cm,height=5cm,
  ylabel={wealth after 300 flips (USD)},ymode=log,
  tick label style={font=\small},label style={font=\small},
  xtick={0,1,2,3},xticklabels={$\tfrac14$ Kelly,$\tfrac12$ Kelly,Kelly,2 Kelly},
  xmin=-0.5,xmax=3.5,ymin=1e-4,ymax=5e6,grid=major,grid style={gray!25},
  legend columns=3,legend style={font=\footnotesize,at={(0.5,-0.25)},anchor=north,draw=none,fill=none,/tikz/every even column/.append style={column sep=8pt}},
  table/col sep=comma]
\addplot[omProp,only marks,mark=triangle*,mark size=2.6pt] table[x=k,y=p90]{figdata/markets-2/29-thinking-in-expected-value/coin.csv};
\addplot[omDef,only marks,mark=*,mark size=2.4pt] table[x=k,y=median]{figdata/markets-2/29-thinking-in-expected-value/coin.csv};
\addplot[omThm,only marks,mark=triangle*,mark options={rotate=180},mark size=2.6pt] table[x=k,y=p10]{figdata/markets-2/29-thinking-in-expected-value/coin.csv};
\addplot[gray,dashed,forget plot] coordinates {(-0.5,25) (3.5,25)};
\legend{90th percentile,median,10th percentile}
\end{axis}
\end{tikzpicture}

\omcaption{Wealth after 300 flips of the 60\% coin from USD~25 (dashed), uncapped, for four
multiples of the Kelly fraction, on a log scale: the 10th, 50th and 90th percentiles of
4\,000 simulated players each. Kelly has the highest median; half Kelly the best
worst tenth; twice Kelly a median below the stake. Data: the chapter's tutorial,
seeded.}\label{fig:m2:thinking-in-expected-value:coin}
\end{omfigure}

\begin{proposition}[Risk of falling by a fraction]\label{prop:m2:thinking-in-expected-value:ruin}
In the continuous-time limit of small bets, staking $c$ times the Kelly fraction, the
probability that wealth ever falls to a fraction $\alpha$ of its starting value is
\[
P = \alpha^{2/c - 1}.
\]
At full Kelly, the chance of ever halving is one half.
\end{proposition}

\begin{proof}
With stake $c\mu/s^2$, log wealth is a Brownian motion with drift $m = (c - \tfrac12
c^2)\mu^2/s^2$ and variance $v = c^2\mu^2/s^2$ per unit time. A Brownian motion with drift
$m > 0$ and variance $v$ ever falls by $d$ with probability $e^{-2md/v}$; with $d = -\ln
\alpha$, $2m/v = 2/c - 1$.
\end{proof}

The formula holds up in simulation (\cref{fig:m2:thinking-in-expected-value:dd}): on a
55\% even-money bet over 4\,000 bets, the share of paths that ever halved was 48.4\% at full
Kelly against 50\% in theory, 11.1\% at half Kelly against 12.5\%. Inverting it gives the
largest multiple of Kelly for a given tolerance: to keep the chance of ever halving
below 10\%, $c = 2/(1 + \ln 0.1/\ln 0.5) = 0.46$. A fall from the peak, which the \omterm{def:m2:thinking-in-expected-value:frac}{drawdown}
measures, recurs over a long enough horizon whatever the fraction: the \omterm{def:m2:thinking-in-expected-value:frac}{risk of ruin} is
the version that has an answer.

\begin{omfigure}
\begin{tikzpicture}
\begin{axis}[width=11cm,height=4.8cm,
  xlabel={multiple of the Kelly fraction, $c$},ylabel={P(wealth ever halves)},
  tick label style={font=\small},label style={font=\small},
  xmin=0.2,xmax=1.55,ymin=0,ymax=0.85,grid=major,grid style={gray!25},
  yticklabel style={/pgf/number format/.cd,fixed,precision=1},
  legend columns=2,legend style={font=\footnotesize,at={(0.5,-0.3)},anchor=north,draw=none,fill=none,/tikz/every even column/.append style={column sep=8pt}},
  table/col sep=comma]
\addplot[omDef,very thick,domain=0.2:1.55,samples=60] {0.5^(2/x - 1)};
\addplot[omThm,only marks,mark=*,mark size=2pt] table[x=c,y=sim]{figdata/markets-2/29-thinking-in-expected-value/drawdown.csv};
\legend{$\alpha^{2/c-1}$ with $\alpha = \tfrac12$,simulated (1\,000 paths)}
\end{axis}
\end{tikzpicture}

\omcaption{Probability that wealth ever falls to half its starting value, against the
multiple of Kelly staked, on a 55\% even-money bet: the continuous-time formula and the
share of 1\,000 simulated paths of 4\,000 bets. Data: the chapter's tutorial,
seeded.}\label{fig:m2:thinking-in-expected-value:dd}
\end{omfigure}

\section{When the edge itself is uncertain}

Everything above assumed the edge known. A trader estimates it, from a backtest or from
the strategy's own history, with an error; and the Kelly fraction is proportional to the
edge, so an overestimated edge means overbetting, which costs more growth than the same
underbetting.

\begin{proposition}[Shrinking for estimation error]\label{prop:m2:thinking-in-expected-value:shrink}
Suppose the mean $\mu$ per unit staked is estimated from $n$ trades by $\hat\mu$, with
standard error $s/\sqrt n$, and the trader stakes $c\hat\mu/s^2$. In the quadratic
approximation the expected \omterm{def:m2:thinking-in-expected-value:growth}{growth rate} is maximised at
\[
c^* = \frac{\mu^2}{\mu^2 + s^2/n} = \frac{n\,\mathrm{SR}^2}{n\,\mathrm{SR}^2 + 1},
\]
where $\mathrm{SR} = \mu/s$ is the Sharpe ratio per trade; $n\,\mathrm{SR}^2$ is the square
of the $t$-statistic of the edge.
\end{proposition}

\begin{proof}
$E[g] = E[f]\mu - \tfrac12 E[f^2]s^2$ with $f = c\hat\mu/s^2$, $E[\hat\mu] = \mu$ and
$E[\hat\mu^2] = \mu^2 + s^2/n$, so $E[g] = c\mu^2/s^2 - \tfrac12 c^2(\mu^2 + s^2/n)/s^2$,
maximised at the stated $c$.
\end{proof}

A strategy whose edge has a $t$-statistic of 1, barely distinguishable from nothing,
should stake half of its estimated Kelly fraction; one with a $t$-statistic of 3, nine
tenths. Staking a fraction of Kelly is, in part, this arithmetic
(\cref{fig:m2:thinking-in-expected-value:uncertain}).

\begin{omfigure}
\begin{tikzpicture}
\begin{axis}[width=11cm,height=4.8cm,
  xlabel={multiple of the estimated Kelly fraction, $c$},ylabel={growth (bp per trade)},
  tick label style={font=\small},label style={font=\small},
  xmin=0,xmax=1.6,ymin=-20,ymax=45,grid=major,grid style={gray!25},
  table/col sep=comma]
\addplot[omDef,very thick,mark=*,mark size=1.4pt] table[x=c,y=g]{figdata/markets-2/29-thinking-in-expected-value/uncertain.csv};
\addplot[omThm,thick,dashed] coordinates {(0.669,-20) (0.669,45)};
\addplot[gray,thin] coordinates {(0,0) (1.6,0)};
\node[font=\footnotesize,anchor=north west] at (axis cs:0.69,-5) {$c^* = 0.669$};
\end{axis}
\end{tikzpicture}

\omcaption{Growth per trade, in basis points, of a trader who estimates the win rate of a
55\% even-money bet from 200 trades and then stakes $c$ times the estimated Kelly
fraction for 500 trades, averaged over 3\,000 repetitions. The best multiple in the
simulation is 0.7; the formula gives 0.669. Staking the full estimated Kelly fraction
earns 19\% less. Data: the chapter's tutorial, seeded.}\label{fig:m2:thinking-in-expected-value:uncertain}
\end{omfigure}

The quadratic approximation also hides fat tails: a trade that can lose more than its
stake, or a return distribution whose extremes are underestimated, calls for smaller
fractions still. And Kelly maximises the growth of one bankroll with no outside
constraints; a fund whose investors leave after a 20\% loss faces a \omterm{def:m2:thinking-in-expected-value:frac}{drawdown}
constraint that Kelly ignores, and should size with the \omterm{def:m2:thinking-in-expected-value:frac}{risk of ruin}, not the \omterm{def:m2:thinking-in-expected-value:growth}{growth
rate}.

\section{Tutorial: repeated bets at several Kelly fractions}\label{tut:m2:thinking-in-expected-value:sizing}

\begin{tutorial}
\textbf{Goal.} Compute the Kelly fraction and \omterm{def:m2:thinking-in-expected-value:growth}{growth rate}, simulate repeated bets at
several multiples of Kelly, check the risk-of-ruin formula, and find the best multiple
when the edge is estimated.
\textbf{End state:} \cref{fig:m2:thinking-in-expected-value:growth,fig:m2:thinking-in-expected-value:coin,fig:m2:thinking-in-expected-value:dd,fig:m2:thinking-in-expected-value:uncertain}
and the numbers of the weekend problem.
\begin{tutsteps}
\item \textbf{Kelly and growth}, and the simulator.
\omcode{code/firm/sizing/firm_sizing.py}{14}{41}{The Kelly fraction, the growth rate and repeated bets.}\label{lst:m2:thinking-in-expected-value:kelly}
\item \textbf{Ruin and estimation error}: the \omterm{def:m2:thinking-in-expected-value:frac}{drawdown} probability, its inverse, and the
shrinkage.
\omcode{code/firm/sizing/firm_sizing.py}{44}{66}{Risk of ruin, the fraction for a tolerance, and shrinkage for an estimated edge.}\label{lst:m2:thinking-in-expected-value:ruin}
\item \textbf{Run} \texttt{sizing\_demo.coin\_table()}, \texttt{sizing\_demo.problem()}
and \texttt{fig\_sizing.py}.
\end{tutsteps}
\textbf{What to change next.} Replay the coin with the 250-dollar cap and compare
constant fractions with a rule that bets just enough to reach the cap; then give the
bets fat-tailed outcomes and find how much the best fraction falls.
\end{tutorial}

\section{Build: the sizing module}\label{bld:m2:thinking-in-expected-value:sizing}

\begin{build}
\bfield{Purpose} Every strategy of the miniature firm is sized here: a fraction of
capital from its estimated edge and risk, cut for estimation error and for the firm's
tolerance of \omterm{def:m2:thinking-in-expected-value:frac}{drawdowns}.
\bfield{Interface} \texttt{kelly\_fraction(p, b)}; \texttt{growth(f, p, b)};
\texttt{simulate(f, p, n, paths, seed, b, cap)}; \texttt{drawdown\_probability(c,
alpha)}; \texttt{fraction\_for\_drawdown(alpha, delta)}; \texttt{shrinkage(n, sharpe)};
\texttt{growth\_uncertain(c, mu, s, n)}.
\bfield{Rules} Wealth multiplicative, stakes a fixed fraction of current wealth; the
\omterm{def:m2:thinking-in-expected-value:frac}{risk of ruin} in its continuous-time form; the shrinkage in the quadratic approximation;
seeded simulations.
\bfield{Acceptance tests} \texttt{code/firm/sizing/tests/}: Kelly maximises growth, twice
Kelly has about zero growth; the ruin formula and its inverse agree; simulated falls to
half match the formula; the shrinkage maximises the expected growth.
\bfield{Stretch} Kelly for several correlated strategies (the fraction becomes
$\Sigma^{-1}\mu$); fat tails; sizing under a hard \omterm{def:m2:thinking-in-expected-value:frac}{drawdown} limit; Bayesian updating of the
edge as trades arrive (One Quant Book 4).
\end{build}

\begin{omsources}
\item V.~Haghani and R.~Dewey, ``Rational decision-making under uncertainty: observed
betting patterns on a biased coin'', October 2016.
\item J.~L.~Kelly, ``A new interpretation of information rate'', Bell System Technical
Journal, 1956.
\end{omsources}

\section{Exercises}

\begin{exercise}[$\star$]\label{exo:m2:thinking-in-expected-value:1}
A bet pays 2 to 1 and wins with probability 0.4. What are its edge and its Kelly fraction?
\end{exercise}

\begin{exercise}[$\star$]\label{exo:m2:thinking-in-expected-value:2}
Why is the expected wealth after 300 flips a poor guide to what a player ends with?
\end{exercise}

\begin{exercise}[$\star$]\label{exo:m2:thinking-in-expected-value:3}
What fraction of the full Kelly \omterm{def:m2:thinking-in-expected-value:growth}{growth rate} does half Kelly keep, in the quadratic
approximation? And a quarter Kelly?
\end{exercise}

\begin{exercise}[$\star\star$]\label{exo:m2:thinking-in-expected-value:4}
At full Kelly, what is the probability of ever falling to 20\% of starting wealth? At half
Kelly?
\end{exercise}

\begin{exercise}[$\star\star$]\label{exo:m2:thinking-in-expected-value:5}
A fund's investors leave if it ever loses 20\% of its starting capital, and it accepts a
10\% chance of that. What multiple of Kelly should it use?
\end{exercise}

\begin{exercise}[$\star\star$]\label{exo:m2:thinking-in-expected-value:6}
An edge is estimated with a $t$-statistic of 2. What multiple of the estimated Kelly
fraction maximises expected growth?
\end{exercise}

\begin{exercise}[$\star\star\star$]\label{exo:m2:thinking-in-expected-value:7}
\emph{Coding.} With \texttt{simulate} and a cap of ten times the stake, estimate the share
of players who reach the cap in 300 flips of the 60\% coin at half Kelly and at Kelly,
and compare with the 95\% quoted by Haghani and Dewey.
\end{exercise}

\begin{exercise}[$\star\star\star$]\label{exo:m2:thinking-in-expected-value:8}
\emph{Find the flaw.} ``Our backtest shows an edge with a Sharpe ratio of 0.1 per trade
over 200 trades, so we stake the full Kelly fraction: it maximises growth.'' Correct it.
\end{exercise}

\section{Problem: The Uncertain Edge}

\begin{problem}[{Weekend problem --- how much of Kelly to stake}]\label{pb:m2:thinking-in-expected-value:1}
A trader has a signal that, it believes, wins 55\% of the time on even-money trades. It
has seen 200 trades, and will size the next 500 as a fraction of capital.

\medskip\noindent\textbf{Part I --- The known edge.}
\begin{enumerate}
\item If the win rate is exactly 55\%, what are the edge, the Kelly fraction and the \omterm{def:m2:thinking-in-expected-value:growth}{growth
rate} per trade?
\item What are the mean and standard deviation of the profit per unit staked, and the
Sharpe ratio per trade?
\item At full Kelly, what is the probability of ever halving the capital?
\item What multiple of Kelly keeps that probability below 10\%?
\item Why does staking twice Kelly destroy the edge?
\end{enumerate}

\medskip\noindent\textbf{Part II --- The estimate.}
\begin{enumerate}[resume]
\item What is the standard error of the win-rate edge estimated from 200 trades, and its
$t$-statistic?
\item What multiple of the estimated Kelly fraction maximises expected growth?
\item What does the simulation give?
\item How much growth does staking the full estimated Kelly fraction lose against the best
multiple?
\item At the best multiple, what is the probability of ever halving?
\end{enumerate}

\medskip\noindent\textbf{Part III --- Robustness.}
\begin{enumerate}[resume]
\item How does the best multiple change with 50 trades of history? With 2\,000?
\item Why is overbetting worse than underbetting by the same amount?
\item What would fat tails do to the best fraction?
\item How would you combine this shrinkage with a \omterm{def:m2:thinking-in-expected-value:frac}{drawdown} limit?
\item What does the coin experiment say about sizing in practice?
\end{enumerate}

\medskip\noindent\textbf{Part IV --- Judgement.}
\begin{enumerate}[resume]
\item Why might a firm size strategies at a quarter Kelly or less?
\item How would you update the fraction as new trades arrive?
\item When is the Kelly framework the wrong one?
\item State the \emph{named result}: the best multiple of Kelly when the edge is estimated
from 200 trades.
\item In one sentence: why stake less than Kelly?
\end{enumerate}
\end{problem}

\section{Interview questions}

\begin{interviewq}[$\star$ \iqroles{trader, researcher}]\label{iq:m2:thinking-in-expected-value:1}
A coin lands heads 60\% of the time and pays even money. How much of your wealth do you
bet, and why?
\end{interviewq}

\begin{interviewq}[$\star$ \iqroles{trader}]\label{iq:m2:thinking-in-expected-value:2}
What is the difference between maximising expected wealth and maximising expected log
wealth?
\end{interviewq}

\begin{interviewq}[$\star\star$ \iqroles{researcher}]\label{iq:m2:thinking-in-expected-value:3}
Derive the Kelly fraction for a continuous return with mean $\mu$ and variance $\sigma^2$.
\end{interviewq}

\begin{interviewq}[$\star\star$ \iqroles{researcher, risk}]\label{iq:m2:thinking-in-expected-value:4}
Why do practitioners use \omterm{def:m2:thinking-in-expected-value:frac}{fractional Kelly}?
\end{interviewq}

\begin{interviewq}[$\star\star$ \iqroles{trader, researcher}]\label{iq:m2:thinking-in-expected-value:5}
You have two uncorrelated strategies with Sharpe ratios 1 and 2. How do you split capital
between them?
\end{interviewq}

\begin{interviewq}[$\star\star\star$ \iqroles{developer}]\label{iq:m2:thinking-in-expected-value:6}
Design a sizing service that sets each strategy's capital daily from its live
performance, with limits.
\end{interviewq}
